\documentclass[../DoAn.tex]{subfiles}
\begin{document}

\section{Mở đầu chương}
Chương này trình bày cơ sở lý thuyết và các phân tích liên quan đến việc xây dựng hệ thống chatbot hỗ trợ giảng viên tự động trả lời câu hỏi của sinh viên qua email. Nội dung chương bao gồm: phân tích bối cảnh và đặc điểm bài toán trong môi trường giáo dục, khảo sát các nghiên cứu và hệ thống tương tự, và giới thiệu các kiến thức nền tảng phục vụ cho việc triển khai hệ thống như: xử lý ngôn ngữ tự nhiên (NLP), mô hình TinyLLaMA, Gmail API và cơ sở dữ liệu PostgreSQL.

\section{Phân tích bài toán và bối cảnh ứng dụng}

\subsection{Thực trạng tương tác qua email giữa sinh viên và giảng viên}
Tại các trường đại học, đặc biệt là các khoa có đông sinh viên, việc sinh viên gửi email hỏi giảng viên là điều diễn ra hàng ngày. Các câu hỏi thường xoay quanh: lịch học, hạn nộp bài, thắc mắc nội dung môn học, thủ tục học vụ,... Đặc điểm chung là nhiều câu hỏi có nội dung giống nhau hoặc tương tự nhau.

Giảng viên phải trả lời nhiều email lặp lại, mất nhiều thời gian, trong khi các email không được trả lời kịp thời có thể ảnh hưởng đến tiến độ học tập và gây khó khăn cho sinh viên. Do đó, có nhu cầu cấp thiết về một hệ thống hỗ trợ trả lời tự động các email có nội dung đơn giản.

\subsection{Mục tiêu và yêu cầu của hệ thống chatbot}
Hệ thống cần giải quyết các yêu cầu sau:
\begin{itemize}
    \item Tự động thu thập email và trích xuất nội dung từ hộp thư giảng viên.
    \item Nhận diện và phân loại email là câu hỏi từ sinh viên hay phản hồi từ giảng viên.
    \item Tạo tập dữ liệu cặp hỏi – đáp để phục vụ huấn luyện chatbot.
    \item Ứng dụng mô hình ngôn ngữ nhẹ (TinyLLaMA) để sinh phản hồi cho các câu hỏi tương tự trong tương lai.
    \item Dễ triển khai, hoạt động tốt trên máy cá nhân hoặc môi trường tài nguyên hạn chế.
\end{itemize}

\section{Các nghiên cứu và hệ thống liên quan}

\subsection{Sự phát triển của hệ thống chatbot}

\subsubsection{Chatbot thế hệ đầu (Rule-based chatbot)}
Các chatbot trước đây chủ yếu hoạt động theo nguyên tắc rule-based (dựa trên quy tắc), sử dụng các kỹ thuật:
\begin{itemize}
    \item \textbf{Pattern matching}: Khớp mẫu từ khóa để xác định ý định người dùng
    \item \textbf{Decision trees}: Cây quyết định để định hướng cuộc hội thoại
    \item \textbf{Template-based responses}: Phản hồi theo mẫu có sẵn
\end{itemize}

Ví dụ điển hình là chatbot ELIZA (1964) sử dụng pattern matching đơn giản để tạo cảm giác hội thoại với người dùng. Hệ thống này hoạt động dựa trên việc nhận diện từ khóa trong câu hỏi và phản hồi theo các mẫu được định sẵn.

\subsubsection{Chatbot thế hệ hiện tại (AI-powered chatbot)}
Ngày nay, các chatbot được phát triển dựa trên:
\begin{itemize}
    \item \textbf{Machine Learning}: Học máy để nhận diện ý định và thực thể
    \item \textbf{Deep Learning}: Mạng nơ-ron sâu cho hiểu ngôn ngữ tự nhiên
    \item \textbf{Large Language Models}: Mô hình ngôn ngữ lớn như GPT, BERT, LLaMA
    \item \textbf{Contextual understanding}: Hiểu ngữ cảnh và duy trì cuộc hội thoại dài
\end{itemize}

\subsection{Các hệ thống hỗ trợ học vụ hiện tại}
Một số hệ thống quản lý học vụ hiện nay như EdStem, Piazza, Moodle Q\&A có hỗ trợ sinh viên đặt câu hỏi và giảng viên trả lời. Tuy nhiên, các hệ thống này chủ yếu yêu cầu truy cập qua nền tảng riêng, không tích hợp với email và không có khả năng tự động trả lời.

Một số chatbot học vụ trên nền tảng Zalo, Facebook Messenger được xây dựng thủ công, với câu trả lời cố định hoặc dựa trên từ khóa, thiếu khả năng học và thích ứng.

\subsection{Các nghiên cứu sử dụng chatbot và dữ liệu email}
Một số nghiên cứu đã thử nghiệm việc dùng chatbot trong môi trường giáo dục, nhưng chủ yếu sử dụng dữ liệu từ diễn đàn, trang hỏi đáp hoặc xây dựng dữ liệu mẫu. Việc sử dụng dữ liệu email thật từ môi trường thực tế (giảng viên và sinh viên) vẫn còn khá hạn chế, đặc biệt là tại Việt Nam.

\subsection{Điểm mới của đề tài}
Khác với các hệ thống hiện có, đề tài này:
\begin{itemize}
    \item Trích xuất dữ liệu trực tiếp từ Gmail của giảng viên.
    \item Tự động phân loại và trích xuất nội dung thành cặp câu hỏi – trả lời từ email.
    \item Huấn luyện mô hình ngôn ngữ trên dữ liệu thật để tạo chatbot có khả năng tự học.
\end{itemize}

\section{Kiến thức nền tảng}

\subsection{Xử lý ngôn ngữ tự nhiên (NLP)}

\subsubsection{Tổng quan về NLP}
Xử lý ngôn ngữ tự nhiên (Natural Language Processing – NLP) là lĩnh vực kết hợp giữa trí tuệ nhân tạo và ngôn ngữ học, nhằm giúp máy tính hiểu và sinh ngôn ngữ như con người. Trong đề tài, NLP được ứng dụng để:
\begin{itemize}
    \item Làm sạch nội dung email (loại bỏ chữ ký, HTML, phần trả lời lặp lại).
    \item Nhận diện và phân loại văn bản thành câu hỏi hoặc câu trả lời.
    \item Chuẩn hóa văn bản đầu vào trước khi huấn luyện mô hình.
\end{itemize}

\subsubsection{Tokenization - Tách từ}
Tokenization là quá trình phân tách văn bản thành các đơn vị nhỏ hơn (tokens) như từ, câu, hoặc ký tự. Đây là bước đầu tiên và quan trọng trong xử lý NLP:
\begin{itemize}
    \item \textbf{Word tokenization}: Tách văn bản thành từng từ
    \item \textbf{Sentence tokenization}: Tách văn bản thành từng câu
    \item \textbf{Subword tokenization}: Tách từ thành các đơn vị nhỏ hơn (BPE, WordPiece)
\end{itemize}

Ví dụ về tokenization:
\begin{verbatim}
Văn bản: "Xin chào, tôi là sinh viên!"
Tokens: ["Xin", "chào", ",", "tôi", "là", "sinh", "viên", "!"]
\end{verbatim}

\subsubsection{Vector Embedding - Biểu diễn vector}
Vector embedding là kỹ thuật chuyển đổi các đơn vị ngôn ngữ (từ, câu, đoạn văn) thành các vector số trong không gian nhiều chiều. Điều này giúp máy tính có thể xử lý và so sánh ý nghĩa ngữ nghĩa:
\begin{itemize}
    \item \textbf{Word embeddings}: Word2Vec, GloVe, FastText
    \item \textbf{Contextualized embeddings}: BERT, RoBERTa, PhoBERT
    \item \textbf{Sentence embeddings}: Sentence-BERT, Universal Sentence Encoder
\end{itemize}

Ví dụ về vector embedding:
\begin{verbatim}
Từ "học sinh" → Vector [0.2, -0.1, 0.8, ..., 0.3] (300 chiều)
Từ "sinh viên" → Vector [0.1, -0.2, 0.7, ..., 0.4] (300 chiều)
\end{verbatim}

Hai vector này sẽ có khoảng cách gần nhau trong không gian vector do có ý nghĩa tương tự.

\subsection{Mô hình TinyLLaMA}
TinyLLaMA là một phiên bản nhỏ gọn của mô hình ngôn ngữ LLaMA do Meta phát triển. Mô hình này có ưu điểm:
\begin{itemize}
    \item Kích thước nhỏ (từ 1B – 3B tham số), phù hợp với máy tính cá nhân.
    \item Dễ huấn luyện lại (fine-tune) trên tập dữ liệu riêng.
    \item Cho kết quả tốt khi huấn luyện trên tập câu hỏi – trả lời đơn giản.
\end{itemize}

\subsection{Gmail API và xác thực OAuth 2.0}
Gmail API là giao diện lập trình của Google cho phép truy xuất, đọc và xử lý email. Các chức năng chính sử dụng trong đề tài:
\begin{itemize}
    \item Truy xuất danh sách chuỗi hội thoại (thread).
    \item Trích xuất nội dung từng email, người gửi, tiêu đề và thời gian.
    \item Áp dụng xác thực OAuth 2.0 để đảm bảo quyền riêng tư và bảo mật.
\end{itemize}

\subsection{PostgreSQL và SQLAlchemy}

\subsubsection{PostgreSQL}
PostgreSQL là hệ quản trị cơ sở dữ liệu mã nguồn mở, mạnh mẽ, được sử dụng để lưu trữ:
\begin{itemize}
    \item Toàn bộ email đã thu thập.
    \item Bảng \texttt{emails}, \texttt{questions}, \texttt{answers} dùng để phân loại và truy vấn dữ liệu.
\end{itemize}

\subsubsection{Cấu trúc dữ liệu}
Hệ thống sử dụng PostgreSQL để lưu trữ email theo dạng câu hỏi và câu trả lời với cấu trúc như sau:

\begin{verbatim}
-- Bảng emails chứa thông tin email gốc
CREATE TABLE emails (
    id SERIAL PRIMARY KEY,
    thread_id VARCHAR(255),
    message_id VARCHAR(255),
    sender VARCHAR(255),
    subject TEXT,
    body TEXT,
    timestamp TIMESTAMP,
    is_question BOOLEAN
);

-- Bảng questions chứa câu hỏi đã được xử lý
CREATE TABLE questions (
    id SERIAL PRIMARY KEY,
    thread_id VARCHAR(255),
    content TEXT,
    processed_content TEXT,
    timestamp TIMESTAMP
);

-- Bảng answers chứa câu trả lời tương ứng
CREATE TABLE answers (
    id SERIAL PRIMARY KEY,
    question_id INTEGER REFERENCES questions(id),
    thread_id VARCHAR(255),
    content TEXT,
    processed_content TEXT,
    timestamp TIMESTAMP
);
\end{verbatim}

\subsubsection{Quy trình nhóm dữ liệu}
Hệ thống lọc theo \texttt{thread\_id} và thời gian gửi email để nhóm lại theo một chuỗi các câu hỏi và câu trả lời liên tục:
\begin{enumerate}
    \item \textbf{Trích xuất email}: Thu thập tất cả email từ Gmail API
    \item \textbf{Phân loại email}: Xác định email nào là câu hỏi, email nào là câu trả lời
    \item \textbf{Nhóm theo thread}: Sử dụng \texttt{thread\_id} để nhóm các email liên quan
    \item \textbf{Sắp xếp theo thời gian}: Sắp xếp các email trong cùng thread theo timestamp
    \item \textbf{Tạo cặp Q\&A}: Ghép câu hỏi với câu trả lời dựa trên thứ tự thời gian
\end{enumerate}

\subsubsection{SQLAlchemy}
SQLAlchemy là công cụ ORM giúp thao tác với cơ sở dữ liệu dễ dàng trong Python.

\subsection{BeautifulSoup và Regex}
Đây là hai công cụ chính để xử lý nội dung email:
\begin{itemize}
    \item \textbf{BeautifulSoup}: Dùng để loại bỏ thẻ HTML khỏi email.
    \item \textbf{Regex (biểu thức chính quy)}: Dùng để tách nội dung hiện tại của email khỏi phần reply cũ, lọc bỏ chữ ký và các đoạn trùng lặp.
\end{itemize}

\section{Xử lý tiền dữ liệu}

\subsection{Thách thức trong xử lý email}
Email trong thực tế thường chứa nhiều thông tin không cần thiết như:
\begin{itemize}
    \item Thẻ HTML từ email định dạng
    \item Phần trích dẫn email cũ (quoted text)
    \item Chữ ký email
    \item Header và metadata
    \item Nội dung forward
    \item Ký tự đặc biệt và markdown
\end{itemize}

\subsection{Các hàm xử lý dữ liệu}

\subsubsection{Hàm extract\_latest\_email\_content}
Hàm này loại bỏ phần reply cũ trong email (hỗ trợ cả tiếng Anh và tiếng Việt) và loại bỏ HTML nếu có:

\textbf{Định nghĩa hàm:}
\begin{verbatim}
def extract_latest_email_content(email_body):
\end{verbatim}

\textbf{Loại bỏ thẻ HTML:}
\begin{verbatim}
    # Loại bỏ tất cả các thẻ HTML
    email_body = re.sub(r'<[^>]+>', '', email_body)
\end{verbatim}

\textbf{Loại bỏ reply cũ:}
\begin{verbatim}
    # Loại bỏ các đoạn reply cũ
    pattern = r'(?i)^((On|Vào) .*?(wrote:|đã viết:))\s*$'
    email_body = re.sub(pattern, '', email_body, 
                       flags=re.MULTILINE)
\end{verbatim}

\textbf{Loại bỏ dòng trích dẫn:}
\begin{verbatim}
    # Loại bỏ các dòng bắt đầu bằng ">"
    email_body = re.sub(r'(?m)^\s*>+.*$', '', email_body)
\end{verbatim}

\textbf{Loại bỏ địa chỉ email:}
\begin{verbatim}
    # Loại bỏ các dòng chứa địa chỉ email
    pattern = r'(?i)^.*?<[a-zA-Z0-9._%+-]+@'
    pattern += r'[a-zA-Z0-9.-]+\.[a-zA-Z]{2,}>.*$'
    email_body = re.sub(pattern, '', email_body, 
                       flags=re.MULTILINE)
\end{verbatim}

\textbf{Xóa khoảng trắng và trả về:}
\begin{verbatim}
    # Xóa các khoảng trắng dư thừa
    email_body = re.sub(r'\n\s*\n+', '\n', 
                       email_body).strip()
    
    return email_body
\end{verbatim}

\subsubsection{Hàm clean\_forwarded\_message}
Hàm này loại bỏ phần forwarded message:

\begin{verbatim}
def clean_forwarded_message(text):
    if not text:
        return text
    split_token = "---------- Forwarded message ---------"
    return text.split(split_token)[0].strip()
\end{verbatim}

\subsubsection{Hàm remove\_html}
Hàm này loại bỏ thẻ HTML và ký tự markdown:

\begin{verbatim}
def remove_html(raw_html):
    # Loại bỏ thẻ HTML
    text = BeautifulSoup(raw_html, "html.parser").get_text()
    # Loại bỏ các dấu '*' hoặc ký tự Markdown
    text = re.sub(r'\*+', '', text)
    return text.strip()
\end{verbatim}

\subsection{Hạn chế và cải tiến}

\subsubsection{Hạn chế hiện tại}
\begin{itemize}
    \item Chưa làm sạch được hoàn toàn dữ liệu
    \item Một số trường hợp đặc biệt chưa được xử lý
    \item Chưa có cơ chế xử lý đa ngôn ngữ hiệu quả
\end{itemize}

\subsubsection{Các cải tiến cần thiết}
\textbf{Xử lý chữ ký email nâng cao:}
\begin{verbatim}
def remove_signature(text):
    signature_patterns = [
        r'Trân trọng.*?$',
        r'Thân ái.*?$',
        r'Best regards.*?$',
        r'Sent from.*?$'
    ]
    for pattern in signature_patterns:
        text = re.sub(pattern, '', text, flags=re.MULTILINE | 
        re.DOTALL)
    return text
\end{verbatim}

\textbf{Chuẩn hóa văn bản:}
\begin{verbatim}
def normalize_text(text):
    # Chuẩn hóa khoảng trắng
    text = re.sub(r'\s+', ' ', text)
    # Xử lý dấu câu
    text = re.sub(r'([.!?])\s*([A-Z])', r'\1 \2', text)
    # Loại bỏ ký tự đặc biệt
    text = re.sub(r'[^\w\s\.\!\?\,\:\;\-\(\)]', '', text)
    return text.strip()
\end{verbatim}

\textbf{Phát hiện ngôn ngữ:}
\begin{verbatim}
def detect_language(text):
    # Sử dụng thư viện langdetect hoặc polyglot
    # để xác định ngôn ngữ và áp dụng xử lý phù hợp
    pass
\end{verbatim}

\section{Kết luận chương}
Chương 2 đã trình bày tổng quan ngữ cảnh bài toán, các nghiên cứu liên quan và nền tảng công nghệ sử dụng trong đề tài. Đặc biệt, chương đã làm rõ:
\begin{itemize}
    \item \textbf{Sự phát triển của chatbot}: Từ rule-based đến AI-powered
    \item \textbf{Các kỹ thuật NLP cơ bản}: Tokenization và vector embedding
    \item \textbf{Kiến trúc dữ liệu}: Cách tổ chức và lưu trữ email trong PostgreSQL
    \item \textbf{Xử lý tiền dữ liệu}: Các phương pháp làm sạch email với những hạn chế cần khắc phục
\end{itemize}

Đây là cơ sở quan trọng để bước sang chương tiếp theo, nơi hệ thống chatbot sẽ được phân tích chi tiết và thiết kế thành các module cụ thể nhằm thực hiện mục tiêu của đồ án.

\end{document}